\section{Discussion}\label{sec:discussion}

\paragraph{What to do differently.}
(1)~Validate patching, steering and KV-cache edits in the format in which the behaviour is used: with the candidates listed after the context, a late patch that changes nearly every free-form answer leaves most listed answers unchanged (\cref{sec:claim3}); localisation methods \citep{mueller2025mib} and an edit's apparent channel (\cref{app:remap}) need the same check.
(2)~With the candidates listed after the context, measure the window per model in that format and edit well inside it (here the late patch from $l_b$, at \NqLbFracMin--\NqLbFracMax{} of the depth, changed only \NqLateListedMin--\NqLateListedMax{} of listed answers; \cref{tab:estimates}); there, a delivered key share well below the natural one predicts a miss. Steering at every position reached named-list answers at every block we tried, but writes at the answer and mostly missed lettered ones (post hoc).
(3)~State the format before reading a layer-wise patching curve, or the answer's attention, at a context token as where the context is read: cutting the answer's attention to the writing token leaves the answer intact with a later list, while without one most answers fail (discovery sample; \cref{app:g-knockout}).

\paragraph{Limitations.}\label{sec:limitations}
(i)~Our passages are counterfactual SQuAD items with added options and the answer named once, scored at the first answer token; naturally occurring later mentions, with the counterfactual written at every one, are the test we most lack, and on passages stage 9 supports only an evaluation warning. (ii)~The confirmation and the window use one story template, and the window three 7--8B models (the trained remap two, with one seed); its depth differs between models, and we claim no universal depth. (iii)~A later mention here repeats the word form (plurals open \GfourteenPluralMin--\GfourteenPluralMax{} of the read; discovery sample); paraphrases and pronouns are untested. (iv)~A list before the writing token gives a small, unexplained negative key read (to $\idKBeforeQwenSeventytwo$ nats at Qwen2.5-72B). (v)~The readers' duplicate-token core depends on how many heads are counted, ablation and knockout disagree on whether the copy route takes over (\cref{app:heads}), and how the answer prefers the moved-to location over the initial one is unresolved. (vi)~The mechanism was tested in two 7B models (and at 1.5B and 3B), the channel-ratio law in three 7--8B models, all in BF16; eager attention, which Gemma-2 requires, lowers the key read by 12--17\,\% in one model (exploratory); the list-after share falls beyond 14B for unknown reasons (discovery sample). % verbatim: "12 % (AFTER) and 17 % (POST)", outcome J part B
(vii)~Values span one to five tokens and prompts are English; the models of stages 8 and 9 are instruction-tuned, and in the base models we ran on IOI a later list opened no positive key read (a weak null; discovery sample, exploratory; \cref{app:g-ioi}). (viii)~We re-tested no third-party channel claim with and without a later mention.

\paragraph{Conclusion.}
Across seven model families on fresh stories, and for an appended options list on counterfactual passages in every model we could score, later mentions decide whether a written value is looked up through its key or copied from its value; on fresh stories the decision reaches generated answers. % seven = \JxFamiliesABCD (design count)
In two 7B models the matching heads write a flag that the answer reads by key and that can move the answer on its own, but is neither specific nor the only carrier. Where the candidates are named also sets when an edit can act: with them listed after the story, patches and single-site edits of the writing token act mainly through its early layers, up to a model-specific depth, and fail silently beyond it; steering at every position escapes with a named list.
